\section{Discussion}
\label{sec:discussion}

\subsection{What the baselines tell us}
\label{sec:disc-insights}

The measured results support four methodological observations.

\paragraph{Simple, well-specified models are strong.} On \task{2}, the official
psychology sheet, a Plackett--Luce model and a kernel-density simulation all
produce identical Kendall $\kendall=0.7673$ in the seed-42 run and
$0.8132\pm0.0342$ across seeds. On \task{3}, the historical \DNF{} rate
($0.3591$) is statistically indistinguishable from the best tree model
($0.3731$, $p=0.749$, negligible effect), and a Bayesian Beta--Binomial
posterior ($0.3725$) is in fact best \emph{on average} across seeds
($0.2938\pm0.0570$ vs $0.2779\pm0.0724$ for the tree model). In both cases a
simple, domain-informed baseline matches or nearly matches more expressive
alternatives. This is not a reason to abandon richer models; it is a calibration
of expectations. A benchmark that reports only ``model A beats model B'' hides
the fact that both may be matched by a good prior at a fraction of the cost.

\paragraph{Generic deep and graph models underperform here.} Two of the six
methodological families do markedly worse than their simpler counterparts, and
the multi-seed results make the gap robust to evaluation-set resampling. On
\task{1}, the LSTM (\maelog{} $=1.0440$, or $0.9436\pm0.0754$ across seeds) is
the \emph{worst} baseline, worse even than the historical mean
($0.6657$), with a large paired effect against the tree model
($\Delta=+1.087$, $p=5.2\times10^{-8}$, Cohen's $d=0.90$). On \task{2}, the GNN
($\kendall=0.5670$, or $0.7440\pm0.1258$ across seeds) trails the psychology
sheet by $0.200$ ($p=3.2\times10^{-7}$, $d=0.56$). Meanwhile rule-aware
features plus tree boosting win \task{1} and are competitive on \task{3}. At
this sample size and feature budget, the domain structure---competitor history,
round format, event identity---appears to matter more than architectural
capacity, and sequence and graph models likely need domain-specific inductive
biases before they can compete. We regard this as the benchmark's clearest
actionable finding.

\paragraph{The hard cases are structural, not incidental.} Errors concentrate
where the data are thin or the phenomenon is stochastic: novices on \task{1}
(\maelog{} $0.4151$), the late test slice on \task{2} ($\kendall=0.6437$) and
\task{3} (\aucpr{} $0.2962$), elites on \task{3} (\aucpr{} $0.1399$), and
cold-start competitors everywhere (\maelog{} $0.9827$). These are not sampling
accidents; they follow from little history, distribution shift and
near-deterministic execution. A useful benchmark must therefore be judged on
the hard subsets, which is why \wcabench{} makes them first-class.

\paragraph{Extrapolation and identification are where claims become fragile.}
\task{4} and \task{5} are included precisely because they are hard to validate.
The \task{4} three-by-three estimates span $3.02$~s (change point), $3.29$~s
(GP+EVT) and $3.74$~s (exponential decay), all far from the community anchor of
$\approx 2.37$~s. None is ``wrong'' in a statistical sense---each is a
consistent fit to its assumptions---but the spread across reasonable methods is
itself the finding, and the aggregate stability scores are not significantly
separated at $21$ events. Similarly, the number of transfer pairs identified
drops from $413$ (correlation) to $312$ (causal forest) to $270$ (IV) to $6$
(DID; $2$ on average across seeds). Reporting the disagreement and the
non-identifiability, rather than a single number, is the scientifically honest
outcome.

\paragraph{Evaluation-set variance matters.} Re-running over three seeds shows
that some conclusions are stable (the tree model wins \task{1} by a wide margin
under every seed) while others are not (the \task{3} winner flips between the
tree model and the Beta--Binomial prior, and the GNN's \task{2} spread,
$\pm0.1258$, is larger than several inter-model gaps). Single-run leaderboards
in this domain should therefore be read with the multi-seed table alongside
them.

\subsection{Towards avoiding saturation}
\label{sec:disc-saturation}

Benchmarks decay: once a task is solved, it stops discriminating. We built two
defences. First, the five tasks form a difficulty spectrum from well-specified
regression to causal identification, so improvements at the easy end do not
exhaust the benchmark. Second, every task carries a hard subset
(Section~\ref{sec:protocol-stratification}); for \task{3}, restricting to
competitors with a historical \DNF{} rate in $[0.1,0.3]$ removes the
easy-to-classify extremes, and for \task{2} the near-tie subset removes rounds
where ranking is trivial. The small-mode results already show that several
subsets are far from saturated, for example elite \task{3} (\aucpr{}
$0.1399$) and cold-start \task{1} (\maelog{} $0.9827$), so the benchmark retains
headroom.

\subsection{Why a domain benchmark, and why speedcubing}
\label{sec:disc-domain}

The coexistence of strong simple models and more expressive alternatives
recalls the long-standing tension between data modelling and algorithmic
prediction \citep{breiman_cultures}: the benchmark's job is to make the
comparison fair, not to anoint a winner. \wcabench{} is a case study in a
broader idea: \emph{rule-governed individual sports} are an underused testbed
for machine learning. Such data combine large scale with genuine
structure---explicit formats, deterministic averaging rules, discrete status
codes---that is invisible to generic pipelines but decisive for accuracy.
Speedcubing is an unusually clean instance: it is fully individual, precisely
timed to the centisecond, globally distributed and publicly released. The
lessons transfer to swimming, track and field, or powerlifting, where analogous
round formats and disqualification rules exist. We therefore see \wcabench{}
less as an endpoint than as a template for building trustworthy benchmarks in
rule-governed athletic domains.

\subsection{Open problems}
\label{sec:disc-open}

Several directions are deliberately left open. Sequence models
(LSTM/Transformer) and graph neural networks currently underperform, so the open
question is not whether they are more accurate in principle but what
domain-specific inductive biases---event-aware embeddings, round-format
constraints, or interaction graphs anchored on head-to-head history---would let
them exploit the structure that tree models already capture. On \task{5}, the
steep drop from $413$ correlational pairs to $270$ IV pairs and $6$ DID pairs
invites more credible causal designs (for example instruments built from the
staggered domestic introduction of events, or synthetic-control designs) while
the difficulty of validating any of them motivates better sensitivity analysis.
Finally, the interaction between non-stationarity and event scale suggests that
per-event hierarchical normalization, rather than a single global
preprocessing, may be the most promising generic improvement.
